\section{Jordan decomposition in a semisimple linear Lie algebra}

\begin{theorem}
    Let $L \subseteq \mathfrak{gl}(V)$ be a semisimple Lie algebra, where $V$ is finite-dimensional.
    Then for every $a \in L$, the semisimple and nilpotent parts of $a$ in $\mathfrak{gl}(V)$ both lie in $L$.
    Hence the abstract Jordan decomposition in $L$ coincides with the usual Jordan decomposition in $\mathfrak{gl}(V)$.
\end{theorem}

\begin{proof}
    Take $a \in L$, and write its usual Jordan decomposition in $\mathfrak{gl}(V)$ as
    \[
        a=a_s+a_n,
        \qquad [a_s,a_n]=0,
    \]
    where $a_s$ is semisimple and $a_n$ is nilpotent.

    Then
    \[
        \operatorname{ad}_{\mathfrak{gl}(V)}(a)
        =
        \operatorname{ad}_{\mathfrak{gl}(V)}(a_s)
        +
        \operatorname{ad}_{\mathfrak{gl}(V)}(a_n)
    \]
    is the Jordan decomposition of $\operatorname{ad}_{\mathfrak{gl}(V)}(a)$ in
    $\mathfrak{gl}(\mathfrak{gl}(V))$. In particular,
    $\operatorname{ad}_{\mathfrak{gl}(V)}(a_s)$ is a polynomial in
    $\operatorname{ad}_{\mathfrak{gl}(V)}(a)$ with no constant term. Since
    $\operatorname{ad}_{\mathfrak{gl}(V)}(a)$ preserves $L$, it follows that
    \[
        [a_s,L]\subseteq L.
    \]
    Thus $\operatorname{ad}(a_s)\vert_L$ is a derivation of $L$.

    Because $L$ is semisimple, every derivation of $L$ is inner. Hence there exists
    $a_s' \in L$ such that
    \[
        [a_s,x]=[a_s',x]
        \qquad \text{for all } x\in L.
    \]
    Therefore
    \[
        [a_s-a_s',L]=0.
    \]

    Now decompose the $L$-module $V$ as
    \[
        V=\bigoplus_i V_i,
    \]
    where each $V_i$ is irreducible. Since $a_s-a_s'$ commutes with the action of $L$,
    Schur's lemma shows that $a_s-a_s'$ acts by a scalar on each $V_i$.

    On the other hand, $L=[L,L]$, so for every $x\in L$ and every $i$,
    \[
        \operatorname{tr}_{V_i}(x)=0.
    \]
    Applying this to $x=a-a_s' \in L$, we get
    \[
        \operatorname{tr}_{V_i}(a-a_s')=0.
    \]
    Since $a_n$ is nilpotent, $\operatorname{tr}_{V_i}(a_n)=0$, hence
    \[
        \operatorname{tr}_{V_i}(a_s-a_s')=0.
    \]
    But $a_s-a_s'$ is scalar on $V_i$, so that scalar must be $0$. Thus
    $a_s-a_s'$ acts trivially on every $V_i$, whence
    \[
        a_s=a_s' \in L.
    \]
    Finally,
    \[
        a_n=a-a_s\in L.
    \]
    So both parts lie in $L$, as required.
\end{proof}

\begin{corollary}
    Let $L$ be a semisimple Lie algebra and $\varphi\colon L\to \mathfrak{gl}(V)$
    a finite-dimensional representation.
    If
    \[
        x=s+n
    \]
    is the abstract Jordan decomposition of $x\in L$, then
    \[
        \varphi(x)=\varphi(s)+\varphi(n)
    \]
    is the usual Jordan decomposition of $\varphi(x)$ in $\mathfrak{gl}(V)$.
\end{corollary}

\begin{proof}
    It is enough to show that
    \[
        \operatorname{ad}_{\varphi(L)}(\varphi(x))
        =
        \operatorname{ad}_{\varphi(L)}(\varphi(s))
        +
        \operatorname{ad}_{\varphi(L)}(\varphi(n))
    \]
    is the abstract Jordan decomposition of
    $\operatorname{ad}_{\varphi(L)}(\varphi(x))$.

    Since $\operatorname{ad}_L(s)$ is semisimple, $L$ is spanned by its eigenvectors.
    Applying $\varphi$, we see that $\varphi(L)$ is spanned by eigenvectors of
    $\operatorname{ad}_{\varphi(L)}(\varphi(s))$. Hence
    $\operatorname{ad}_{\varphi(L)}(\varphi(s))$ is semisimple.

    Since $\operatorname{ad}_L(n)$ is nilpotent, its induced action on the quotient
    $\varphi(L) \cong L / \ker(\varphi)$ is nilpotent. Hence
    $\operatorname{ad}_{\varphi(L)}(\varphi(n))$ is nilpotent.

    Also, because $[s,n]=0$,
    \[
        \bigl[\operatorname{ad}_{\varphi(L)}(\varphi(s)),
            \operatorname{ad}_{\varphi(L)}(\varphi(n))\bigr]=0.
    \]
    Thus the displayed decomposition is the abstract Jordan decomposition in the
    semisimple linear Lie algebra $\varphi(L)$. By the theorem above, it is also the
    usual Jordan decomposition in $\mathfrak{gl}(V)$.
\end{proof}

